% ============================================================
% appendix.tex - 附录
% ============================================================

\appendix
\section*{附录}
\addcontentsline{toc}{section}{附录}

% ==============================
\subsection*{A. 源代码}
% ==============================

% 代码完整原版粘贴，使用 lstlisting 格式。
% 按问题编号分节。

\subsubsection*{A.1 问题一求解代码}

\begin{lstlisting}[language=Python, caption={问题一求解代码 (question1\_solve.py)}]
# -*- coding: utf-8 -*-
"""
问题一求解代码
"""
import numpy as np
# ... 完整代码粘贴于此 ...
\end{lstlisting}

\subsubsection*{A.2 问题二求解代码}

\begin{lstlisting}[language=Python, caption={问题二求解代码 (question2\_solve.py)}]
# ... 完整代码粘贴于此 ...
\end{lstlisting}

% ==============================
\subsection*{B. 运行数据与补充表格}
% ==============================

% 正文中未展示的详细数据表格、中间计算结果等。
% 使用 longtable 处理大表格。

\begin{longtable}{cccc}
  \caption{补充数据表} \\
  \toprule
  \textbf{序号} & \textbf{数据1} & \textbf{数据2} & \textbf{数据3} \\
  \midrule
  \endfirsthead
  \caption{补充数据表（续）} \\
  \toprule
  \textbf{序号} & \textbf{数据1} & \textbf{数据2} & \textbf{数据3} \\
  \midrule
  \endhead
  \bottomrule
  \endfoot

  1 & -- & -- & -- \\
  % ...
\end{longtable}
